% !TEX program = xelatex
% code: numbered blocks, kinds, \listofcodes, \begin{code}, lines= and
% \codedir, output and terminal, \iosample, inline \code in arguments,
% Agda / Datalog / Stata / Haskell / SQL
\documentclass{akizumi-note}
\akizumiNewCodeKind[env=example, counter=code]{example}{Example}{範例}
\akizumiNewCodeKind[env=recipe]{recipe}{Recipe}{作法}
\codedir{build/}
\begin{filecontents*}[overwrite]{code-demo.py}
import sys

def add(a, b):
    return a + b  # line 4

print(add(1, 2))
\end{filecontents*}
\begin{filecontents*}[overwrite]{code-run.txt}
$ python code-demo.py
3
\end{filecontents*}
\title{Code}
\begin{document}
\maketitle
\listofcodes

\section{The \code{a\_b} function and \code[cpp]{std::vector<int>}}

Inline: \code{x_1 % #{} \n} and \code|a{b| and \code[python]{print("hi")  # note}.
Bold: \textbf{\code[python]{print("\n")}}; in a footnote\footnote{See \code{f(x)_2}.}.

\begin{theorem}[\code{sort} is stable]\label{thm:s}
  Merge sort with \code[cpp]{a[i] <= b[j]} keeps equal keys in order.
\end{theorem}

\begin{codeblock}[numbered, label=lst:fit, title=fit.py, lang=py]
def fit(x, y):
    return (x.T @ x).inverse() @ x.T @ y   # least squares
\end{codeblock}

\begin{codeblock}[lang=Python, title=unnumbered.py]
pass
\end{codeblock}

\begin{example}{python}{(sum.py) Adding numbers}
total = sum([1, 2, 3])
\end{example}

\begin{code}[label=lst:idiom]{python}{Loop over a list}
for x in xs:
    print(x)
\end{code}

\begin{recipe}[numbers=false]{haskell}{Fold}
sum' :: [Int] -> Int
sum' = foldr (+) 0  -- right fold
\end{recipe}

\cref{lst:fit} and \cref{lst:idiom} (numbers shared with examples), \cref{thm:s}.

\inputcode[lang=Python, lines=3-4, title=code-demo.py (lines 3--4)]{code-demo.py}
\inputcode[output, title=output]{code-run.txt}
\inputcode[numbered, label=lst:whole]{code-demo.py}

\begin{terminal}[title=zsh]
$ ax build
>>> 1 + 1
2
\end{terminal}

\iosample{1 2 \\ 3 4}{[3, 7] \\ done}
\iosample{}{[89.2, 74.8] \\}

\begin{codeblock}[lang=Agda, title=Nat.agda]
data ℕ : Set where
  zero : ℕ
  suc  : ℕ → ℕ  -- successor
{- block
   comment -}
\end{codeblock}

\begin{codeblock}[lang=Datalog]
# reachability
path(X, Y) :- edge(X, Y).
path(X, Z) :- path(X, Y), edge(Y, Z), not blocked(Z).
\end{codeblock}

\begin{codeblock}[lang=Stata]
* first stage
use "supas.dta", clear
reg yeduc high##young if young == 1, r   // robust
\end{codeblock}

\begin{codeblock}[lang=SQL]
SELECT name FROM pokemon WHERE hp > 150; -- strong
\end{codeblock}

\begin{verbatim}
plain verbatim stays as it was
\end{verbatim}
\end{document}
